% PTE Core 模考 1B —— 口语 错题与详解·整理卷  版本 002（单文件自包含）
% 编译: xelatex 本文件.tex   （需 Noto CJK + TeX Gyre Heros 字体；无外部图片）
\documentclass[11pt]{article}

% --------- 样式（原 ptestyle.sty，已内联）---------
% ============================================================
%  ptestyle.sty  —  PTE 模考错题整理卷 共享样式
%  编译引擎: XeLaTeX  (中英混排, 需要 Noto CJK 字体)
% ============================================================

\usepackage[a4paper,margin=2cm,headsep=4mm]{geometry}
\usepackage{fontspec}
\usepackage{xeCJK}
\usepackage[table]{xcolor}
\usepackage[normalem]{ulem}        % \sout 删除线
\usepackage{tabularx}
\usepackage{array}
\usepackage{ragged2e}
\usepackage{enumitem}
\usepackage[most]{tcolorbox}
\usepackage{fancyhdr}
\usepackage{graphicx}
\usepackage{needspace}
\usepackage{tikz}

% ---------- 字体 ----------
\setmainfont{TeX Gyre Heros}            % 拉丁: Helvetica 风格(纤细清晰)
\setCJKmainfont{Noto Sans CJK SC}
\setCJKsansfont{Noto Sans CJK SC}
\newfontfamily\cjkserif{Noto Serif CJK SC}
\linespread{1.18}
\setlength{\parindent}{0pt}
\setlength{\parskip}{4pt}

% ---------- 配色 ----------
\definecolor{cWrong}{HTML}{D32F2F}     % 红  = 修改 / 错误
\definecolor{cFix}{HTML}{1B7F3B}       % 绿  = 正确建议
\definecolor{cDelete}{HTML}{1565C0}    % 蓝  = 删除
\definecolor{cInsert}{HTML}{7B1FA2}    % 紫  = 插入
\definecolor{cPartial}{HTML}{E07B00}   % 橙  = 部分得分
\definecolor{cNote}{HTML}{0277BD}      % 蓝  = 注解
\definecolor{cAccent}{HTML}{16A085}    % 主题青绿(平台主色)
\definecolor{cWrongBg}{HTML}{FBE0E0}
\definecolor{cDelBg}{HTML}{DCE8FB}
\definecolor{cInsBg}{HTML}{EEE0F5}
\definecolor{cHead}{HTML}{20303A}
\definecolor{cGrayBox}{HTML}{F4F6F7}
\definecolor{cModelBg}{HTML}{EAF7EF}
\definecolor{cYouBg}{HTML}{FFF7F5}
\definecolor{cNoteBg}{HTML}{EAF3FA}
\definecolor{cRule}{HTML}{D7DEE2}

% ---------- 行内批改宏 (复刻平台标注) ----------
% \fix{原词}{正确}  红色删除线 + 绿色上标改正  (修改/拼写/空格)
\newcommand{\fix}[2]{%
  \colorbox{cWrongBg}{\textcolor{cWrong}{\sout{#1}}}%
  \,\textsuperscript{\textcolor{cFix}{\footnotesize #2}}}
% \del{原词}  蓝色删除线  (应删除)
\newcommand{\del}[1]{%
  \colorbox{cDelBg}{\textcolor{cDelete}{\sout{#1}}}%
  \,\textsuperscript{\textcolor{cDelete}{\footnotesize 删}}}
% \ins{词}  紫色下划线  (应插入)
\newcommand{\ins}[1]{%
  \textsuperscript{\textcolor{cInsert}{\footnotesize $\wedge$}}%
  \colorbox{cInsBg}{\textcolor{cInsert}{#1}}}
% \miss{词}  灰色删除线  (口语漏读 / 未作答)
\newcommand{\miss}[1]{\textcolor{gray}{\sout{#1}}}
% \add{原词}{改正}  橙色波浪线 + "补"标签  (平台漏标、本卷补标的错误)
\definecolor{cAdd}{HTML}{E65100}
\newcommand{\add}[2]{%
  \textcolor{cAdd}{\uwave{#1}}%
  \,\textsuperscript{\textcolor{cAdd}{\footnotesize 补：#2}}}
% 高亮(口语发音): 好/中/差
\newcommand{\spk}[2]{\textcolor{#1}{#2}}

% ---------- 评分单元格着色 ----------
\newcommand{\sfull}[1]{\textcolor{cFix}{\bfseries #1}}     % 满分
\newcommand{\spart}[1]{\textcolor{cPartial}{\bfseries #1}} % 部分
\newcommand{\szero}[1]{\textcolor{cWrong}{\bfseries #1}}   % 零/低分

% ---------- 分数小标签 ----------
\newtcbox{\chip}[1][cAccent]{on line, boxsep=2pt, left=4pt, right=4pt, top=1pt,
  bottom=1pt, arc=3pt, colback=#1, colframe=#1, coltext=white, nobeforeafter,
  fontupper=\bfseries\footnotesize}

% ---------- 题块小标题 ----------
\newcommand{\ptesub}[1]{%
  \par\needspace{2\baselineskip}\smallskip
  {\color{cAccent}\rule[-0.2ex]{3pt}{1.05em}}\hspace{6pt}%
  {\bfseries\color{cHead}#1}\par\nopagebreak\smallskip}

% ---------- 题块外框 ----------
% \begin{qblock}{标号·题型}{右上信息(得分/用时)} ... \end{qblock}
\newtcolorbox{qblock}[2]{
  breakable, enhanced, arc=2.5pt,
  colback=white, colframe=cRule, boxrule=0.6pt,
  left=11pt, right=11pt, top=8pt, bottom=10pt,
  toptitle=3pt, bottomtitle=3pt, lefttitle=11pt, righttitle=11pt,
  colbacktitle=cHead, coltitle=white, fonttitle=\bfseries,
  title={#1\hfill\normalfont\bfseries\small #2}}

% ---------- 文本块: 题干 / 范文 / 你的作答 / 建议 ----------
\newtcolorbox{promptbox}{enhanced, breakable, colback=cGrayBox, colframe=cGrayBox,
  arc=2pt, left=8pt, right=8pt, top=6pt, bottom=6pt, boxrule=0pt}
\newtcolorbox{modelbox}{enhanced, breakable, colback=cModelBg, colframe=cModelBg,
  borderline west={3pt}{0pt}{cFix}, arc=1pt, left=8pt, right=8pt, top=6pt, bottom=6pt, boxrule=0pt}
\newtcolorbox{youbox}{enhanced, breakable, colback=cYouBg, colframe=cYouBg,
  borderline west={3pt}{0pt}{cWrong}, arc=1pt, left=8pt, right=8pt, top=6pt, bottom=6pt, boxrule=0pt}
\newtcolorbox{notebox}{enhanced, breakable, colback=cNoteBg, colframe=cNoteBg,
  borderline west={3pt}{0pt}{cNote}, arc=1pt, left=8pt, right=8pt, top=6pt, bottom=6pt, boxrule=0pt}

% ---------- 评分表 ----------
% 用法见正文; 这里给表头宏
\newcommand{\scorehead}{%
  \rowcolors{2}{cGrayBox}{white}
  \arrayrulecolor{cRule}}

% ---------- 章节大标题 ----------
\newcommand{\ptesection}[3]{% #1=中文名 #2=英文名 #3=该项分数
  \needspace{6\baselineskip}
  \begin{tcolorbox}[enhanced, colback=cAccent, colframe=cAccent, sharp corners,
     left=14pt, right=14pt, top=10pt, bottom=10pt, boxrule=0pt]
    {\fontsize{20}{24}\selectfont\bfseries\color{white} #1\ \ }%
    {\large\color{white!85} #2}\hfill
    {\large\color{white}本项得分\ \ }{\fontsize{22}{24}\selectfont\bfseries\color{white} #3}%
    \ {\color{white!85}/ 90}
  \end{tcolorbox}}

% ---------- 题型小节分隔 (口语 RA/RS/DI/RTS/ASQ 等) ----------
\newcommand{\ptetask}[2]{%
  \needspace{4\baselineskip}\medskip
  \begin{tcolorbox}[enhanced, colback=cHead, colframe=cHead, sharp corners,
     left=12pt, right=12pt, top=5pt, bottom=5pt, boxrule=0pt]
    {\large\bfseries\color{white} #1}\ \ {\color{white!80} #2}
  \end{tcolorbox}\smallskip}

% ---------- 颜色图例 ----------
\newcommand{\ptelegend}{%
  \begin{tcolorbox}[enhanced, colback=cGrayBox, colframe=cRule, boxrule=0.5pt,
    arc=2pt, left=10pt, right=10pt, top=6pt, bottom=6pt]
  {\bfseries\color{cHead}颜色说明\quad}
  \fix{改}{正}~修改/拼写\quad
  \del{删}~应删除\quad
  \ins{插}~应插入\quad
  \add{补}{改正}~平台漏标(本卷补标)\quad
  {\color{cFix}\rule{1em}{0.6em}}~范文/正确\quad
  {\color{cNote}\rule{1em}{0.6em}}~AI 建议
  \end{tcolorbox}}

% ---------- 页眉页脚 ----------
\pagestyle{fancy}
\fancyhf{}
\renewcommand{\headrulewidth}{0.4pt}
\renewcommand{\footrulewidth}{0pt}
\fancyhead[L]{\footnotesize\color{cHead}\bfseries PTE Core 模考 1B}
\fancyhead[R]{\footnotesize\color{gray} 错题与详解整理卷}
\fancyfoot[C]{\footnotesize\color{gray}\thepage}

\begin{document}

% --------- 封面（原 cover.tex）---------
% ---------- 封面 / 成绩总览 ----------
\definecolor{mListen}{HTML}{1F3A5F}
\definecolor{mRead}{HTML}{C2D70B}
\definecolor{mSpeak}{HTML}{4D4D4D}
\definecolor{mWrite}{HTML}{A01B7D}

\begin{titlepage}
\thispagestyle{empty}
\centering
\vspace*{1.4cm}

{\fontsize{30}{36}\selectfont\bfseries\color{cHead} PTE Core 模考 1B}\\[4mm]
{\fontsize{17}{20}\selectfont\color{cAccent}\bfseries 错题与详解 · 整理卷}\\[3mm]
{\color{gray} APEUni / 猩际 PTE 模考　·　提交时间 2026-06-21 12:36}\\[1.2cm]

% ---- 总分 ----
\begin{tcolorbox}[width=0.62\textwidth, halign=center, enhanced,
  colback=cAccent, colframe=cAccent, arc=6pt, boxrule=0pt, top=10pt, bottom=10pt]
  {\color{white}\large 总分 Overall}\\[2pt]
  {\fontsize{46}{50}\selectfont\bfseries\color{white} 40}
\end{tcolorbox}

\vspace{1.0cm}

% ---- 四项分数条形图 ----
\begin{tcolorbox}[width=0.84\textwidth, enhanced, colback=white, colframe=cRule,
  boxrule=0.8pt, arc=4pt, left=18pt, right=18pt, top=14pt, bottom=14pt]
\setlength{\tabcolsep}{6pt}\renewcommand{\arraystretch}{1.7}
\sffamily
\begin{tabular}{@{}r@{\hskip 8pt}l@{\hskip 10pt}l@{}}
 \bfseries Listening 听力 & {\color{mListen}\rule{3.67cm}{0.42cm}} & {\bfseries\large 33}\\
 \bfseries Reading 阅读   & {\color{mRead}\rule{4.67cm}{0.42cm}}   & {\bfseries\large 42}\\
 \bfseries Speaking 口语  & {\color{mSpeak}\rule{3.22cm}{0.42cm}}  & {\bfseries\large 29}\\
 \bfseries Writing 写作   & {\color{mWrite}\rule{6.33cm}{0.42cm}}  & {\bfseries\large 57}\\
\end{tabular}\\[2mm]
{\footnotesize\color{gray} 满分 90　|　刻度: 条长 = 得分 / 90}
\end{tcolorbox}

\vfill
\begin{flushleft}\small\color{gray}
\textbf{说明:}本卷由本次模考的题目与"详解"截图逐题誊写、重新排版而成,完整保留每道题的
\emph{题干 / 原文、范文、你的作答、逐项评分、彩色批改与 AI 中文建议}。\\
所有错误按平台标注用颜色标出(见下方图例)。原始"详解"PDF 不可复制、仅截图,本卷内容均为人工对照誊录,若与原图有出入以原图为准。
\end{flushleft}
\vspace{4mm}
\ptelegend
\end{titlepage}

% --------- 口语正文（原 sec_speaking_body.tex）---------
% ===================== 口语 Speaking =====================
\ptesection{口语}{Speaking · RA + RS + DI + RTS + ASQ}{29}

\vspace{1mm}
{\small\color{gray} 本部分依据「口语详解1」逐题誊录，含 \textbf{RA 朗读}$\times$7、\textbf{RS 复述句子}$\times$11、
\textbf{DI 描述图}$\times$3、\textbf{RTS 情景应答}$\times$3、\textbf{ASQ 简答}$\times$1。每题保留平台 AI 评分（各维度）与你的语音识别结果。
失分高度集中在 \textbf{发音 Pronunciation} 与 \textbf{流利度 Fluency}：内容/恰当性多为中上，但发音与流利度普遍偏低，是本项仅得 29 的主因。}
\vspace{1mm}

\begin{notebox}
{\footnotesize 说明：朗读/复述题段落为按截图誊录的题面文字（已去除停顿斜线与口误重复）；
DI、RTS、ASQ 的「你的作答」为平台 \textbf{AI 语音识别}原文（含识别/发音错误，原样保留，以反映真实作答）。
分数标注：满分/高分 {\color{cFix}绿}、中等 {\color{cPartial}橙}、低分 {\color{cWrong}红}；各维度满分见表头。}
\end{notebox}

% ============================================================
\ptetask{RA 朗读}{Read Aloud · 内容 Content / 5　发音 Pronunciation / 90　流利度 Fluency / 90}

{\small\color{gray} 7 题内容均 \sfull{5/5}（题面读全），失分全在发音与流利度。}\smallskip

\begin{notebox}
{\footnotesize\textbf{平台 AI 共性建议（反复出现）：}①你朗读得不够流畅，有较多卡顿或迟疑；
②你的弱读单词（如 and / but / or / it / is / that）发音没有弱化到位；
③你单词发音过慢或有拖音现象，建议读快一点点；④读错的单词偏多，重点关注标红/标黄的词后重读纠正。}
\end{notebox}

\begin{qblock}{RA 1\quad \#608\ \ Summerhill School}{总分 \textcolor{cAccent}{\bfseries 12} / 90}
\begin{promptbox}
Summerhill School was regarded with considerable suspicion by the educational establishment. Lessons were optional for pupils at the school, and the government of the school was carried out by a School Council, of which all the pupils and staff were members, with everyone having equal voting rights.
\end{promptbox}
{\small 评分：内容 \sfull{5/5}\quad·\quad 发音 \spart{14/90}\quad·\quad 流利度 \szero{10/90}}
\end{qblock}

\begin{qblock}{RA 2\quad Statistical chances (universities)}{总分 \textcolor{cAccent}{\bfseries 18} / 90}
\begin{promptbox}
The survey found that the statistical chances of someone from a poor background being accepted at one of the country's most respected universities are far lower than those of a student from a wealthy family. This means that the inequalities in society are likely to be passed down from one generation to the next.
\end{promptbox}
{\small 评分：内容 \sfull{5/5}\quad·\quad 发音 \spart{22/90}\quad·\quad 流利度 \szero{15/90}}
\end{qblock}

\begin{qblock}{RA 3\quad Language diversity}{总分 \textcolor{cAccent}{\bfseries 16} / 90}
\begin{promptbox}
Despite a number of events in recent years devoted to language diversity, language endangerment and multilingualism, such as the International Year of Languages, public awareness of the issues is still remarkably limited. Only one in four of the population know that half the languages of the world are so seriously endangered that they are unlikely to survive the present century.
\end{promptbox}
{\small 评分：内容 \sfull{5/5}\quad·\quad 发音 \spart{20/90}\quad·\quad 流利度 \szero{11/90}}
\end{qblock}

\begin{qblock}{RA 4\quad European \& American orchestras}{总分 \textcolor{cAccent}{\bfseries 18} / 90}
\begin{promptbox}
The advantage of the great European and American orchestras is that they were able to establish their iconic status in an age when their identity could become entrenched. There was less competition and it was easier to create a brand. Not only did they have the best halls, they attracted the best musicians, who tended to stay put.
\end{promptbox}
{\small 评分：内容 \sfull{5/5}\quad·\quad 发音 \spart{23/90}\quad·\quad 流利度 \szero{13/90}}
\end{qblock}

\begin{qblock}{RA 5\quad Loggerhead turtle}{总分 \textcolor{cAccent}{\bfseries 16} / 90}
\begin{promptbox}
It is time for this young loggerhead turtle to go to work. We can tether turtles in these little cloth harnesses, put them onto this tank-like swimming place. University of North Carolina biologist Ken Lohmann studies sea turtles that are programmed from birth for an extraordinary journey. Mother turtles nest the eggs on the beach and then return to the sea, and the eggs hatch about 50 to 60 days later.
\end{promptbox}
{\small 评分：内容 \sfull{5/5}\quad·\quad 发音 \spart{22/90}\quad·\quad 流利度 \szero{10/90}}
\end{qblock}

\begin{qblock}{RA 6\quad Agricultural science}{总分 \textcolor{cAccent}{\bfseries 17} / 90}
\begin{promptbox}
While advances in agricultural science have always been critical to ensuring we help feed the world, their impact and importance is even greater now as population grows at a rapid rate and the availability of arable land steadily declines. Science and technology solutions are essential to meeting growing demand for food, maintaining market competitiveness, and adapting to and mitigating risks.
\end{promptbox}
{\small 评分：内容 \sfull{5/5}\quad·\quad 发音 \spart{21/90}\quad·\quad 流利度 \szero{12/90}}
\end{qblock}

\begin{qblock}{RA 7\quad Globalization}{总分 \textcolor{cAccent}{\bfseries 12} / 90}
\begin{promptbox}
The benefits and disadvantages of globalization are the subject of ongoing debate. The downside to globalization can be seen in the increased risk for the transmission of diseases. Globalization has, of course, led to great good too. Richer nations can now come to the aid of poorer nations in crisis. Increasing diversity in many countries has meant more opportunity to learn about and celebrate other cultures.
\end{promptbox}
{\small 评分：内容 \sfull{5/5}\quad·\quad 发音 \spart{15/90}\quad·\quad 流利度 \szero{10/90}}
\end{qblock}

% ============================================================
\ptetask{RS 复述句子}{Repeat Sentence · 内容 Content / 3　发音 / 90　流利度 / 90}

{\small\color{gray} 11 题。下表为题面句子与各维度得分；内容分越低表示漏词/错词越多。}\smallskip

{\small\scorehead
\begin{tabularx}{\linewidth}{@{}c >{\RaggedRight\arraybackslash}X c c c c@{}}
\rowcolor{cHead}
\textcolor{white}{\bfseries\#} & \textcolor{white}{\bfseries 题面句子 Sentence} &
\textcolor{white}{\bfseries 总分} & \textcolor{white}{\bfseries 内容/3} &
\textcolor{white}{\bfseries 发音} & \textcolor{white}{\bfseries 流利度}\\
1428 & Today's lecture is canceled because the lecturer is ill. & \textcolor{cAccent}{\bfseries 48} & 2.2 & 59 & 59\\
1430 & I am available this Thursday afternoon. & \textcolor{cAccent}{\bfseries 32} & 1.5 & 64 & 54\\
1449 & Tuition fees will vary according to the field of study. & \textcolor{cAccent}{\bfseries 34} & 1.6 & 55 & 61\\
1431 & Newspapers around the country are reporting the stories of the president. & \textcolor{cAccent}{\bfseries 24} & 1.2 & 59 & 49\\
1429 & The department determines whether or not the candidates pass. & \textcolor{cWrong}{\bfseries 10} & 0.8 & 37 & 21\\
1487 & A very basic feature of computing is counting and calculating. & \textcolor{cAccent}{\bfseries 27} & 1.7 & 51 & 37\\
1450 & Note that the deadline of the submission of proposals has been extended for a week. & \textcolor{cAccent}{\bfseries 35} & 1.8 & 52 & 50\\
1451 & All students must return the books to the college library before the end of the term. & \textcolor{cAccent}{\bfseries 28} & 1.4 & 47 & 57\\
1448 & The current labor force is more competitive than it has been for a long time. & \textcolor{cAccent}{\bfseries 42} & 1.9 & 56 & 64\\
1432 & Foods containing overabundant calories supply little or no nutritional value. & \textcolor{cWrong}{\bfseries 13} & 0.6 & 51 & 61\\
1488 & Contemporary literature works have been broadened and extended through interpretation. & \textcolor{cWrong}{\bfseries 10} & 0.2 & 48 & 38\\
\end{tabularx}}
\smallskip
{\footnotesize\color{gray} 注：\#1429 / \#1432 / \#1488 内容分极低（0.2$\sim$0.8 / 3），多处漏词或识别错误，是 RS 部分主要失分点。}

% ============================================================
\ptetask{DI 描述图}{Describe Image · 内容 Content / 5　发音 / 90　流利度 / 90}

{\small\color{gray} 3 题。内容（要点覆盖）多为 4.5$\sim$4.7 / 5，但发音、流利度仍偏低。「你的作答」为 AI 语音识别原文。}\smallskip

\begin{qblock}{DI 1\quad \#577\ \ 人口年龄分布图}{总分 \textcolor{cAccent}{\bfseries 57} / 90}
\ptesub{图表 Image}
\begin{promptbox}
柱状图 + 饼图：按年龄段（5--14、15--24、25--34、35--44、45--54、55--64、65--74）显示人口占比（Per cent），25--34 段占比最高（约 16.2\%）。
\end{promptbox}
\ptesub{你的作答 Your Response（AI 语音识别）}
\begin{youbox}
ok look at these pictures we can see a x is present and the white is agent group and as below the left is may lands a right is female and we can see that from twenty five to thirty four is the will sits a large people and i drewn i five eight choose a thirty four for the female is the most part of them in like eighteen five
\end{youbox}
{\small 评分：内容 \sfull{4.7/5}\quad·\quad 发音 \spart{71/90}\quad·\quad 流利度 \szero{33/90}}
\end{qblock}

\begin{qblock}{DI 2\quad \#567\ \ 澳大利亚手机品牌（2009 \& 2019）}{总分 \textcolor{cAccent}{\bfseries 55} / 90}
\ptesub{图表 Image}
\begin{promptbox}
对比饼图：``Popular Mobile Phone Brands in Australia, Year 2009 \& 2019''，比较两个年份各手机品牌（如 Apple 等）的市场占比变化。
\end{promptbox}
\ptesub{你的作答 Your Response（AI 语音识别）}
\begin{youbox}
ok look at this chart is about popular mobile phone brands in australia year in nineteen nine and nineteen nineteen and twenty nineteen a we can see that ins a left is about twenty and nine is the most popular mobile phone brand is apple and from the left is a twenty nineteen the most popular's mobile phone is still is essential and the next year april
\end{youbox}
{\small 评分：内容 \sfull{4.5/5}\quad·\quad 发音 \spart{73/90}\quad·\quad 流利度 \szero{35/90}}
\end{qblock}

\begin{qblock}{DI 3\quad \#576\ \ 欧洲地图}{总分 \textcolor{cAccent}{\bfseries 48} / 90}
\ptesub{图表 Image}
\begin{promptbox}
地图：欧洲（提及 Australia、Romania、Germany 等地名——作答中地名多有口误），标示若干国家/区域位置。
\end{promptbox}
\ptesub{你的作答 Your Response（AI 语音识别）}
\begin{youbox}
look at this image a we can say it's a map of about european a like australia and romania and g choice and we can see the span is no a deft blows a maps and we can see is a feet and is a sins um a bit slow and a force of the map and the germany and australia is in the center of the map
\end{youbox}
{\small 评分：内容 \sfull{4.5/5}\quad·\quad 发音 \spart{64/90}\quad·\quad 流利度 \szero{28/90}}
\end{qblock}

% ============================================================
\ptetask{RTS 情景应答}{Respond to a Situation · 恰当性 Appropriateness / 90　发音 / 90　流利度 / 90}

{\small\color{gray} 3 题（PTE Core 题型）。给出情景，需口头回应。「你的作答」为 AI 语音识别原文。}\smallskip

\begin{qblock}{RTS 1\quad \#20\ \ Barbecue Party}{总分 \textcolor{cAccent}{\bfseries 24} / 90}
\ptesub{情景 Situation}
\begin{promptbox}
You've organized a barbecue with your friends at your place. However, upon their arrival, you notice some of your friends are using your kitchen utensils carelessly. You need to remind everyone to handle the utensils with care to maintain their quality and keep ourselves safe. What would you say to your friends?
\end{promptbox}
\ptesub{你的作答 Your Response（AI 语音识别）}
\begin{youbox}
hi i notice that some of my friends are using a my kitchen uses careless i think i need to remind everyone to handle a announces with care to maintain your quality and keep ourselves safe so i should say that you should handle your own immunities to care about your own qualities and keep ourselves safe
\end{youbox}
{\small 评分：恰当性 \spart{40/90}\quad·\quad 发音 \szero{21/90}\quad·\quad 流利度 \szero{19/90}}
\end{qblock}

\begin{qblock}{RTS 2\quad \#21\ \ Team-building Activity}{总分 \textcolor{cAccent}{\bfseries 19} / 90}
\ptesub{情景 Situation}
\begin{promptbox}
You've organized a team-building activity for your colleagues at a local park. However, upon arrival, you notice some team members are not participating actively and seem disengaged. You need to remind everyone of the importance of teamwork and encourage active participation to maximize the benefits of the activity. What would you say to your colleagues?
\end{promptbox}
\ptesub{你的作答 Your Response（AI 语音识别）}
\begin{youbox}
hi i think i have noticed that sound him members are not participating activity and seems disease juice i think i need to remind everyone if you to the importance of teamwork and encourage active participant to maximize the benefit of the i so i should i should say that every one of our team should know the importance of the teamwork and you should be active um be outgoing that's all
\end{youbox}
{\small 评分：恰当性 \spart{34/90}\quad·\quad 发音 \szero{14/90}\quad·\quad 流利度 \szero{15/90}\quad{\footnotesize\color{gray}（答题用时 01:30）}}
\end{qblock}

\begin{qblock}{RTS 3\quad \#22\ \ Community Garden Project}{总分 \textcolor{cAccent}{\bfseries 29} / 90}
\ptesub{情景 Situation}
\begin{promptbox}
{\color{cWrong}\footnotesize[本题原题文字在截图中被详解弹窗遮挡/分页处未截全，未能誊录；仅保留题号与你的作答。主题为「社区花园 Community Garden Project」。]}
\end{promptbox}
\ptesub{你的作答 Your Response（AI 语音识别）}
\begin{youbox}
hi i have notice that some gardens fields are neglected who growing rapits i think i need to remind everyone of their responsibility to maintain their asylum cities to maintain the success of the garden so i should say that do other things to a members ms science aimed to ensure the success of their gardens
\end{youbox}
{\small 评分：恰当性 \spart{44/90}\quad·\quad 发音 \szero{25/90}\quad·\quad 流利度 \szero{24/90}}
\end{qblock}

% ============================================================
\ptetask{ASQ 简答}{Answer Short Question}

\begin{qblock}{ASQ\quad \#1189}{——}
\ptesub{问题 Question}
\begin{promptbox}
If someone tells you the truth, what is the opposite?
\end{promptbox}
\ptesub{参考答案 Answer}
\begin{modelbox}
Falsity\ /\ falseness\ /\ untruth
\end{modelbox}
{\footnotesize\color{gray} 你的作答：已录音（约 8 秒），截图未提供识别文字与单题得分。详解中仅含此 1 题 ASQ。}
\end{qblock}

\end{document}
